\documentclass[11pt,a4paper]{article}

% ========================
% ESSENTIAL PACKAGES
% ========================

% Language and encoding
\usepackage[utf8]{inputenc}
\usepackage[T1]{fontenc}
\usepackage[english]{babel}
%\usepackage{natbib}

% Math packages
\usepackage{amsmath}        % Enhanced math environments
\usepackage{amssymb}        % Additional math symbols
\usepackage{amsthm}         % Theorem environments
\usepackage{amsfonts}       % Additional math fonts
\usepackage{mathtools}      % Extension to amsmath
\usepackage{bm}             % Bold math symbols
\usepackage{dsfont}         % Double-struck fonts (for sets like ℝ, ℕ)
\usepackage{mathrsfs}       % Script fonts for math

% Graphics and figures
\usepackage{graphicx}       % Include graphics
\usepackage{float}          % Better float placement
\usepackage{subcaption}     % Subfigures
\usepackage{tikz}           % Drawing diagrams
\usepackage{pgfplots}       % Plotting
\pgfplotsset{compat=1.18}

% Page layout and formatting
\usepackage[margin=1in]{geometry}
\usepackage{fancyhdr}       % Custom headers and footers
\usepackage{titlesec}       % Customize section titles
\usepackage{enumitem}       % Customize lists
\usepackage{parskip}        % Paragraph spacing

% Colors and highlighting
\usepackage{xcolor}
\usepackage{mdframed}       % Framed environments
\usepackage{tcolorbox}      % Colored boxes

% Tables
\usepackage{array}          % Extended array environments
\usepackage{booktabs}       % Professional tables
\usepackage{multirow}       % Multi-row cells

% References and links
\usepackage{hyperref}       % Hyperlinks
\usepackage{cleveref}       % Smart cross-references

% Additional useful packages
\usepackage{cancel}         % Cancel terms in equations
\usepackage{siunitx}        % SI units
\usepackage{listings}       % Code listings (if needed)

\usepackage{authblk}

\usepackage[authoryear]{natbib}

% ========================
% CUSTOM COMMANDS
% ========================

% Common math sets
\newcommand{\R}{\mathbb{R}}
\newcommand{\N}{\mathbb{N}}
\newcommand{\Z}{\mathbb{Z}}
\newcommand{\Q}{\mathbb{Q}}
\newcommand{\C}{\mathbb{C}}

% Common operators
\DeclareMathOperator{\lcm}{lcm}
\DeclareMathOperator{\Tr}{Tr}
\DeclareMathOperator{\rank}{rank}
%\DeclareMathOperator{\span}{span}
%\DeclareMathOperator{\null}{null}

% Shortcuts for common expressions
\newcommand{\abs}[1]{\left|#1\right|}
\newcommand{\norm}[1]{\left\|#1\right\|}
\newcommand{\inner}[2]{\left\langle #1, #2 \right\rangle}
\newcommand{\ceil}[1]{\left\lceil #1 \right\rceil}
\newcommand{\floor}[1]{\left\lfloor #1 \right\rfloor}

% Derivatives and integrals
\newcommand{\dd}[1]{\,\mathrm{d}#1}
\newcommand{\dv}[2]{\frac{\mathrm{d}#1}{\mathrm{d}#2}}
\newcommand{\pdv}[2]{\frac{\partial #1}{\partial #2}}

\renewcommand{\hat}{\widehat}
\renewcommand{\tilde}{\widetilde}
\renewcommand{\bar}{\overline}

\newcommand{\vct}{\boldsymbol}
\newcommand{\tr}{\mathrm{tr}}

% ========================
% THEOREM ENVIRONMENTS
% ========================

\theoremstyle{definition}
\newtheorem{definition}{Definition}[section]
\newtheorem{example}{Example}[section]
\newtheorem{exercise}{Exercise}[section]

\theoremstyle{plain}
\newtheorem{theorem}{Theorem}[section]
\newtheorem{lemma}{Lemma}[section]
\newtheorem{corollary}{Corollary}[section]
\newtheorem{proposition}{Proposition}[section]
\newtheorem{assumption}{Assumptio}[section]

\theoremstyle{remark}
\newtheorem{remark}{Remark}[section]
\newtheorem{note}{Note}[section]

\newcommand{\supp}{\mathrm{supp}}
\newcommand{\op}{\mathrm{op}}

\newcommand{\mX}{\mathcal X}
\newcommand{\mD}{\mathcal D}
\newcommand{\bP}{\mathbb P}
\newcommand{\lp}{L^2(\bP)}
\newcommand{\lx}{L^2(\mX)}
\newcommand{\ltwo}{L^2(\mX)}
\newcommand{\ud}{\mathrm d}
\newcommand{\mH}{\mathcal H}
\newcommand{\mA}{\mathcal A}

\newcommand{\diag}{\mathrm{diag}}

\newcommand{\mI}{\mathcal I}

% ========================
% COLORED BOXES
% ========================

% Important theorem box
\newtcolorbox{importantbox}{
	colback=blue!5!white,
	colframe=blue!75!black,
	title=Important
}

% Warning/Note box
\newtcolorbox{notebox}{
	colback=yellow!10!white,
	colframe=orange!75!black,
	title=Note
}

% Example box
\newtcolorbox{examplebox}{
	colback=green!5!white,
	colframe=green!75!black,
	title=Example
}

% ========================
% DOCUMENT SETUP
% ========================

% Page style
\pagestyle{fancy}
\fancyhf{}
\rhead{\thepage}
\lhead{\leftmark}


% Hyperlink setup
\hypersetup{
	colorlinks=true,
	linkcolor=blue,
	filecolor=magenta,      
	urlcolor=cyan,
	citecolor=red
}

\title{Honest and Adaptive Pointwise Confidence Intervals in Linear Functional Regression: An Anisotropic Kernel Ridge Regression Approach}
\author[1]{Paromita Dubey \footnote{Author names listed in alphabetical order.}}
\author[2]{Zhiyuan Tang}
\author[2]{Yining Wang}
\affil[1]{Marshall Business School, University of Southern California}
\affil[2]{Naveen Jindal School of Management, University of Texas at Dallas}


\begin{document}

\maketitle

\begin{abstract}
Abstract here.
\end{abstract}

\section{Introduction}

We consider the classical model of linear functional regression with data $\{X_i,Y_i\}_{i=1}^n$, where $X_1,\cdots,X_n\sim P_X$ are independently and identically distributed \emph{functions} 
whose domain is a non-empty compact interval $I\subseteq\mathbb R$, and $\{Y_i\}_{i=1}^n$ are scalar response variables, modeled by
\begin{equation}
\mathbb E[Y_i|X_i] = \eta_0(X_i) := \gamma_0 + \int_I X_i(t)\beta_0(t)\ud t,
\label{eq:model}
\end{equation}
where $\gamma_0\in\mathbb R$ is the intercept term and $\beta_0: I\to\mathbb R$ is the linear functional, both being unknown parameters.
Many well-known methods have been proposed to estimate $\gamma_0$ and $\beta_0$ from data, such as functional principal component analysis (fPCA) approaches \citep[Chapter 8]{ramsay2005functional},
\citep{cai2006prediction} and methods based on reproducing kernel Hilbert spaces \citep{yuan2010reproducing,cai2012minimax}.

In this paper, we consider the question of constructing \emph{pointwise} confidence intervals on $\eta_0(x)$ which are valuable in quantifying the uncertainty in estimates of the functional model
when it is applied to a future data curve $x$. We study the construction of \emph{honest} confidence intervals, which are functions $\varphi$ mapping a data curve $x:I\to\mathbb R$ to a compact interval $\varphi(x)\subseteq\mathbb R$
such that
\begin{equation}
\liminf_{n\to\infty}\inf_{\eta_0\in\Theta} \Pr\left[\eta_0(x)\in\varphi(x)\right] \geq 1-\alpha,\;\;\;\;\;\;\forall x\in\mathrm{supp}(P_X).
\label{eq:honest-ci}
\end{equation}
where $\alpha\in(0,1)$ is a pre-determined significance level,
$\supp(P_X)$ is the support of $P_X$, $\Theta$ is a fixed and known function class and the randomness in Eq.~(\ref{eq:honest-ci}) is over $\{(X_i,Y_i)\}_{i=1}^n$.
To understand how \emph{adative} an honest confidence interval is, we further study the worst-case mean-squared widths of constructed intervals
\begin{equation}
\sup_{\eta_0\in\Theta}\mathbb E_{\eta_0} \left[\int_{\supp(P_X)}\big|\varphi(x)\big|^2\ud P_X(x)\right],
\label{eq:adaptive}
\end{equation}
where $|\varphi(x)|$ is the Lesbesgue measure of the output interval $\varphi(x)$. We are particularly interested in how the above worst-case risk converges to zero as $n\to\infty$,
and how it compares with the mean-squared prediction errors
$$
\sup_{\eta_0\in\Theta}\mathbb E_{\eta_0} \left[\int_{\supp(P_X)}\int_I [x(t)(\hat\eta_n(t)-\eta_0(t)]^2\ud t\ud P_X(x)\right]
$$
for classical fPCA or RKHS estimators $\hat\eta$ of $\eta_0$.

\subsection{Our contributions}

\subsection{Related works}

\subsection{Notations}

For two sequences of real numbers $\{a_n\}$ and $\{b_n\}$, we write $a_n\lesssim b_n$, or equivalently $a_n=O(b_n)$, if there exists a numerical constant $C$ such that $|a_n|\leq C|b_n|$ for all $n$.
We write $a_n\gtrsim b_n$ or $a_n=\Omega(b_n)$ if $b_n\lesssim a_n$, and $a_n\asymp b_n$ or $a_n=\Theta(b_n)$ if both $a_n\lesssim b_n$ and $a_n\gtrsim b_n$ hold.
We denote $L_2(I):=\{x: I\to\mathbb R: \int_I x(t)^2\ud t<+\infty\}$ as the class of all square integrable functions on $I$.
For two functions $x,z\in L_2(I)$, $\langle x,z\rangle$ denotes the standard inner product of $L_2(I)$; that is, $\langle x,z\rangle = \int_I x(t)z(t)\ud t$.

\section{Functional regression via kernel ridge regression}

In this section, we give an overview of an reproducing kernel Hilbert space (RKHS) approach towards linear functional regression, following the seminal works of \cite{cai2012minimax} which also lays the foundation of 
methodological and theoretical contributions of this paper.
To clarify, all results in this section are relatively well-known in the literature, and we are summarizing them for the purpose of completeness only.

Let $K:I\times I\to\mathbb R$ be a real, symmetric, square integrable, and non-negative definite function, known as the \emph{kernel function}.
There then exists a reproducing kernel Hilbert space $\mH\subseteq L_2(I)$ associated with inner product $\langle\cdot,\cdot\rangle_{\mH}$
such that the following reproducing property holds for any $f\in\mH$:
$$
f(t) = \langle K(t,\cdot), f\rangle_{\mH}\;\;\;\;\;\;\forall t\in I.
$$
Assuming $\beta_0\in\mH$. Then a standard kernel ridge regression (KRR) approach to estimate $\alpha_0,\beta_0$ from $\{X_i,Y_i\}_{i=1}^n$ is to solve the following penalized optimization problem
\begin{equation}
\hat\gamma_n,\hat\beta_n \in \arg\min_{\gamma\in\mathbb R,\beta\in\mH} \left\{\frac{1}{n}\sum_{i=1}^n (Y_i-\gamma-\langle X_i,\beta\rangle)^2 + \lambda_n \|\beta\|_\mH^2\right\},
\label{eq:krr}
\end{equation}
where $\lambda_n>0$ is a penalization parameter.
Let $\bar X=\frac{1}{n}\sum_{i=1}^n X_i$ and $\bar Y=\frac{1}{n}\sum_{i=1}^n Y_i$ be the sample averages. It is easy to verify that $\hat\gamma_n=\bar Y$ and 
\begin{equation}
\hat\beta_n\in\arg\min_{\beta\in\mH}\left\{\frac{1}{n}\sum_{i=1}^n (Y_i'-\langle X_i',\beta\rangle)^2 + \lambda_n\|\beta\|_\mH^2\right\}
\label{eq:krr-nointercept}
\end{equation}
where $Y_i'=Y_i-\bar Y$ and $X_i'=X_i-\bar X$.

\subsection{Reduction to an optimization problem in $L_2(I)$}

The work of \cite{cai2012minimax} makes the observation that Eq.~(\ref{eq:krr-nointercept}) can be equivalently reformulated into an infinite-dimensional optimization problem
involving the standard inner product and norm of $L_2(I)$ only, which is helpful for implementation and theoretical analysis purposes.
We summarize here how such equivalence is built below.

By Mercer theorem, there exists an orthonormal basis $\{\phi_j\}_{j=1}^{\infty}$ of $L_2(I)$
and a sequence of non-increasing, strictly positive real numbers $\rho_1\geq\rho_2\geq\cdots$ such that $\phi_j\in\mH$ for all $j$, and furthermore
$$
\langle\phi_j,\phi_k\rangle = \int_I \phi_j(t)\phi_k(t)\ud t = \delta_{jk},\;\;\;\;\;\langle\phi_j,\phi_k\rangle_{\mH}= \frac{\delta_{jk}}{\rho_j},\;\;\;\;\;\;\int_I K(t,\cdot)\phi_j(t)\ud t=\rho_j\phi_j.
$$
where $\delta_{jk}=\vct 1\{j=k\}$ is the Kronecker delta notation.
Note that because $\mH$ is known, $\{\phi_j,\rho_j\}_{j=1}^{\infty}$ are also known.
Following notations in \citep{cai2012minimax}, we define linear operator $L_{K^{1/2}}:L^2(I)\to L^2(I)$ and its inverse operator $L_{K^{-1/2}}$ as follows:
$$
L_{K^{1/2}}(\phi_j) = \sqrt{\rho_j}\phi_j,\;\;\;\;\;\; L_{K^{-1/2}}(\phi_j) = \frac{1}{\sqrt{\rho_j}}\phi_j,\;\;\;\;\;\forall j\geq 1.
$$
In other words, for any $f=\sum_{j=1}^{\infty}\alpha_j\phi_j\in L^2(I)$ such that $\sum_j\alpha_j^2<+\infty$, $L_{K^{1/2}}(f)=\sum_{j=1}^{\infty}\alpha_j\sqrt{\rho_j}\phi_j$
and therefore $L_{K^{1/2}}f\in\mH$ because $\sum_{j=1}^{\infty}\langle L_{K^{1/2}}f,\phi_j\rangle^2/\rho_j <+\infty$.
In fact, \citep[Proposition 2]{yuan2010reproducing} shows that for $f\neq 0$ and $L_{K^{1/2}}f\neq 0$, $f\in L^2(I)$ if and only if $L_{K^{1/2}}f\in\mH$ and furthermore $\int_I f(t)^2\ud t=\|f\|_2^2=\|L_{K^{1/2}}f\|_\mH^2$.
The optimal solution $\hat\beta_n$ in Eq.~(\ref{eq:krr-nointercept}) can then be expressed as $\hat\beta_n = L_{K^{1/2}}\hat\theta_n$, where
\begin{equation}
\hat\theta_n \in \arg\min_{\theta\in L_2(I)}\left\{\frac{1}{n}\sum_{i=1}^n (Y_i'-\langle L_{K^{1/2}}X_i', \theta\rangle^2 + \lambda_n\|\theta\|_2^2\right\},
\label{eq:krr-nointercept-l2}
\end{equation}
where $\|\theta\|_2^2 =\int_I \theta(t)^2\ud t$ is the standard squared $L_2$ norm on $I$.
For brevity, we will also use the notation of $Z_i' = L_{K^{1/2}}X_i'$ throughout this paper.

\subsection{Kernel ridge regression via spectral decomposition}

The infinite-dimensional optimization problem in Eq.~(\ref{eq:krr-nointercept-l2}) can be reduced to a finite-dimensional problem and efficiently solved
using Gram matrices, known as the kernel trick. In this section, we show an alternative characterization of $\hat\theta_n$, representing it as a smoothed version of functional PCA
which is useful motivating our anisotropic estimators later.

Let
$$
\hat T_n := \frac{1}{n}\sum_{i=1}^n Z_i'\otimes Z_i'
$$
be the sample covariance operator, where $Z_i'=L_{K^{1/2}}X_i'$. By the spectral theorem, there exists orthonormal $\{\hat\psi_j\}_{j=1}^n$ in $L_2(I)$ and $\hat s_1\geq \hat s_2\geq\cdots\geq \hat s_n\geq 0$ such that
$$
\hat T_n = \sum_{j=1}^n \hat s_j\hat \psi_j\otimes\hat \psi_j.
$$
The estimator $\hat\theta_n$ can then be expressed as
\begin{equation}
\hat\theta_n = (\hat T_n+\lambda_n I)^{-1}\frac{1}{n}\sum_{i=1}^n Y_i'Z_i' = \frac{1}{n}\sum_{i=1}^n Y_i'\sum_{j=1}^{\infty}\frac{\hat\nu_j(Z_i')}{\hat s_j+\lambda}\hat\psi_j,
\label{eq:krr-nointercept-fpca}
\end{equation}
where $\hat\nu_j(Z_i') = \langle Z_i', \hat\psi_j\rangle=\int_I Z_i'(t)\hat\psi_j(t)\ud t$.

\section{Assumptions}

In this section we list and discuss assumptions that we impose throughout this paper. For some results, additional assumptions are needed: 
we will remark on them when such additional assumptions are introduced specifically for selected theorems or lemmas.

We first introduce mild regularity conditions to make sure the problem is well-posed.
\begin{enumerate}
\item[(A1)] Noise variables $\varepsilon_i:=Y_i-\eta_0(X_i)$ are centered, independent, homoskedastic with variance $\sigma_\varepsilon^2<+\infty$, and satisfies $\Pr[|\varepsilon_i|>t]\leq 2e^{-C_\varepsilon t^2/2}$
for some $C_\varepsilon<+\infty$ and all $t>0$;
\item[(A2)] $\beta_0\in\mH$ and $\|\beta_0\|_\mH\leq 1$, which is equivalent to $\|\theta_0\|_2=\|L_{K^{-1/2}}\beta_0\|_2 =(\int_I \theta_0(t)^2\ud t)^{1/2} \leq 1$.
\item[(A3)] For $x\sim P_X$, $\|L_{K^{1/2}}x\|_2= (\int_I (L_{K^{1/2}}x)(t))^2\ud t)^{1/2}\leq 1$ almost surely. We also assume that $x\neq 0$ and $L_{K^{1/2}}x\neq 0$ almost surely, to avoid degenerate situations.
\end{enumerate}

We remark that in Assumptions (A2) and (A3) we use one as the upper bounds to $\|\beta_0\|_\mH$ and $\|x\|_2$ for convenience only. Our constructed honest confidence bands do not require such knowledge
and they remain valid with 1 being replaced by any other numerical constants.

We use $\Theta_2(\mH)$ to denote the class of all regression functionals $\beta_0$ that satisfy (A2). Note that $\Theta_2(\mH)$ depends on the selection of the RKHS $\mH$, but not on the covariate distribution $P_X$.

Let $Z:=L_{K^{1/2}}X$ be a random function and $P_Z$ be the distribution of $Z$, which is induced by both $P_X$ and $K$. Consider the population covariance operator
$$
T := \mathbb E_{P_Z}\left[Z\otimes Z\right].
$$
By the spectral theorem, there exists an orthonormal basis $\{\psi_j\}_{j=1}^{\infty}$ of $L_2(I)$ and $s_1\geq s_2\geq\cdots\geq 0$ such that
\begin{equation}
T = \sum_{j=1}^{\infty} s_j\psi_j\otimes\psi_j.
\label{eq:specT}
\end{equation}
Note that, in general, because the population covariance operator $T$ is unknown, $\{\psi_j,s_j\}$ are also unknown.

The following assumptions on imposed directly on the structure of $\{\psi_j\}_j$ and $\{s_j\}_j$.
\begin{enumerate}
\item[(B1)] There exists numerical constants $r>0$ and $0<c_s\leq C_s<+\infty$ such that $c_s j^{-2r}\leq s_j\leq C_sj^{-2r}$ for all $j\geq 1$;
\item[(B2)] There exists a numerical constant $C_p<+\infty$ such that for every $j,k\geq 1$, $\mathbb E[\langle Z,\psi_j\rangle^2\langle Z,\psi_k\rangle^2] \leq C_p\times \mathbb E[\langle Z,\psi_j\rangle^2]\times \mathbb E[\langle Z,\psi_k\rangle^2]= C_ps_js_k$;
\item[(B3)] There exists numerical constants $\kappa\in[0,1/4)$ and $C_\nu<+\infty$ such that for every $j\geq 1$, $|\langle Z,\psi_j\rangle| \leq \sqrt{s_j}\cdot n^{\kappa}$ almost surely.
\end{enumerate}

We remark that, in Assumption (B3), the smaller $\kappa$ the stronger the assumption is. For the strongest version of $\kappa=0$, it essentially assumes $|\langle Z,\psi_j\rangle|\lesssim j^{-r}$,
which coincides with coefficient decay assumptions adopted in many fPCA analysis.
{\color{red}In particular, the results in \cite{cai2006prediction} shows that in this case ($\kappa=0$ in (B3)) fPCA achieves an MSE rate of $n^{-(2r-1)/(2r+1)}$ assuming no decay of $|\langle\theta_0,\psi_j\rangle|$.}

\section{Confidence intervals via isotropic under-smoothing}

As a warm-up exercise and a backup method with minimal assumptions, we first consider honest confidence intervals constructed by \emph{under-smoothing} an isotropic kernel ridge estimator.
Let $\{\lambda_n\}_{n=1}^{\infty}$ be a sequence of positive penalty parameters that satisfies
\begin{equation}
\lambda_n n\log^2 n \to 0,\;\; \lambda_n n^{1-\delta}\to \infty\;\;\;\;\;\;\text{as}\;\;n\to\infty
\label{eq:lambdan-1}
\end{equation}
for any $\delta>0$.
Let $\hat\gamma_n = \bar Y$ and $\hat\beta_n=L_{K^{1/2}}\hat\theta_n$ where 
$$
\hat\theta_n = (\hat T_n+\lambda_n I)^{-1}\frac{1}{n}\sum_{i=1}^n Y_i'Z_i',
$$
with $Y_i'=Y_i-\bar Y$, $Z_i'=L_{K^{1/2}}(X_i-\bar X)$ and $\hat T_n=\frac{1}{n}\sum_{i=1}^n Z_i'\otimes Z_i'$.
This is an \emph{under-smoothed} kernel ridge regression estimator with $\lambda_n=o(n^{-1})$, while the optimal penalty parameter choice for estimating $\eta_0$ in mean-squared errors under $P_X$
is $\lambda_n\asymp n^{-2r/(2r+1)}$, per results in \cite{cai2012minimax}.

For a test function $x$ in the support $P_X$, an honest and adaptive confidence interval of $\eta(x)$ at significance level $\alpha\in(0,1)$ is constructed as
$$
\varphi(x) = [\hat\eta_n(x)-\tau_{\alpha/2}\hat\sigma_n(L_{K^{1/2}}x),\;\; \hat\eta_n(x)+\tau_{\alpha/2}\hat\sigma_n(L_{K^{1/2}}x)]
$$
where $\tau_{\alpha/2}$ is the $(1-\alpha/2)$-quantile of the standard Normal distribution, and
$$
\hat\eta_n(x) = \hat\gamma_n + \int_I x(t)\hat\beta_0(t)\ud t,\;\;\;\;\;\;\hat\sigma_n(z) = \frac{\sigma_\varepsilon}{\sqrt{n}}\sqrt{\langle z, (\hat T_n+\lambda_n I)^{-1}z\rangle}.
$$
%To facilitate our analysis, we also introduce the following quantity $\tilde z_n(z)$ defined as
%$$
%\tilde\sigma_n(z) := \frac{\sigma_\varepsilon}{\sqrt{n}}\sqrt{\langle z, (\hat T_n+\lambda_n I)^{-1}\hat T_n(\hat T_n+\lambda_n I)^{-1}z\rangle},
%$$
%which is smaller than $\hat\sigma_n(z)$ almost surely.

\subsection{Validity and rates of convergence}

\begin{theorem}
Suppose Assumptions (A1), (A2), (B1), (B2), (B3) hold where (B3) holds for some $\kappa<1$. Suppose also that Eq.~(\ref{eq:lambdan-1}) holds. Then 
for any $x\in\supp(P_X)$, conditioned on the event that
\begin{equation}
\liminf_{n\to\infty}\inf_{\beta_0\in\Theta_2(\mH)} \Pr[\eta_0(x)\in\varphi(x)] \geq 1-\alpha\;\;\;\;\;\;\forall x\in\supp(P_X).
\label{eq:simple-validity}
\end{equation}
Furthermore, as $n\to\infty$, 
\begin{equation}
\sup_{\beta_0\in\Theta_2(\mH)}\mathbb E_{\eta_0}\left[\int_{\supp(P_X)}\big|\varphi(x)\big|^2\ud P_X(x)\right]\lesssim n^{-\frac{2r-1}{2r}+\delta}
\label{eq:simple-validity}
\end{equation}
for any $\delta>0$.
\label{thm:simple}
\end{theorem}

\begin{remark}
If $\{\varepsilon_i\}_{i=1}^n$ are Normally distributed random variables, then Eq.~(\ref{eq:simple-validity}) in Theorem \ref{thm:simple}) holds without assuming (B3).
\end{remark}


\subsection{Minimax optimality}

\begin{theorem}
Fix arbitrary $\alpha\in(0,1)$, $r>0$ and $\kappa=0$. There exists a distribution $P_X$ satisfying Assumptions (A1), (A3), (B1), (B2), and (B3), such that for any $\varphi$ satisfying Eq.~(\ref{eq:honest-ci})
with respect to the constructed $P_X$ and $\Theta_2(\mH)$, it holds that
$$
\sup_{\beta_0\in\Theta_2(\mH)} \mathbb E_{\beta_0}\left[\int_{\supp(P_X)}\big|\varphi(x)\big|^2\ud P_X(x)\right] \gtrsim n^{-\frac{2r-1}{2r}}.
$$
\label{thm:simple-lb}
\end{theorem}

\section{Confidence intervals via anisotropic under-smoothing}

Let
$$
\hat\theta_n = (\hat T_n+\Lambda_n)^{-1}\frac{1}{n}\sum_{i=1}^n Y_i' Z_i'
$$
where $\Lambda_n$ is a self-adjoint covariance operator of the form
$$
\Lambda_n = \sum_{j=1}^n \lambda_{nj}\hat\psi_j\otimes\hat\psi_j
$$
for some $\{\lambda_{nj}\}_{j=1}^n\geq 0$, and the entire $\Lambda_n$ depends only on $\{Z_i'\}_{i=1}^n$ but not $Y_i'$.
For the purpose of this section, we adopt the choice of 
$$
\lambda_{nj} = \sqrt{\hat s_j\lambda_n}.
$$

{\color{blue}
Intuition: the squared bias term, after Cauchy-Schwarz, is
\begin{align*}
\sum_{j=1}^n \hat\nu_j(z)^2\frac{\lambda_{nj}^2}{(s_j+\lambda_{nj})^2} = \sum_{j=1}^n \lambda_n \frac{\hat\nu_j(z)^2}{(\sqrt{\hat s_j}+\sqrt{\lambda_n})^2}.
\end{align*}

the variance term is
\begin{align*}
\frac{1}{n}\sum_{i=1}^n\langle z, (\hat T_n+\Lambda_n)^{-1}Z_i\rangle^2 &= \langle z, (\hat T_n+\Lambda_n)^{-1}\hat T_n(\hat T_n+\Lambda_n)^{-1}z\rangle\\
&= \sum_{j=1}^n \hat\nu_j(z)^2\frac{\hat s_j}{(\hat s_j+\lambda_{nj})^2} = \sum_{j=1}^n \frac{\hat\nu_j(z)^2}{(\sqrt{\hat s_j}+\sqrt{\lambda_n})^2}.
\end{align*}
}

\section{More adaptive intervals with anisotropic KRR}


\section{Proofs}

For brevity we assume $\mathbb E[x]=0$, $\gamma_0=0$ and $X_i'=X_i$ throughout all proofs, which is conventional in the analysis of linear functional regression in the literature \citep{cai2012minimax,yuan2010reproducing}.
For all \emph{lower bound} proofs, we also assume $L_{K^{1/2}}=L_{K^{-1/2}}$ is the identity operator to simplify the proofs.



\subsection{Some important technical lemmas}

We first establish a few technical lemmas which are important to our proof of later results.
\begin{lemma}
Suppose Assumptions (A1) and (A3) hold. Suppose Assumption (B3) holds for some $\kappa<1/2$, and $\lambda_n$ satisfies Eq.~(\ref{eq:lambdan-1}). Then there exists a constant $n_0\in\mathbb N$ depending only on $\kappa<1/2$ such that for all $n\geq n_0$, with probability at least $1-n^{-1}$,
$$
\frac{1}{2}(T+\lambda_n I)\preceq \hat T_n+\lambda_n I\preceq 2(T+\lambda_n I)
$$
where $\hat T_n=\frac{1}{n}\sum_{i=1}^n Z_i\otimes Z_i$, $T=\mathbb E[Z\otimes Z]$, $Z_i=L_{K^{1/2}}X_i$, $Z=L_{K^{1/2}}Z$, and for two self-adjoint operators $A,B$, $A\preceq B$ if $\langle x,Ax\rangle\leq \langle x, Bx\rangle$
for any square integrable on function $x:I\to\mathbb R$.
\label{lem:empirical-population-cov}
\end{lemma}

We remark that similar versions of Lemma \ref{lem:empirical-population-cov} have appeared, in different forms, in classical analysis of kernel ridge regression estimators and RKHS approaches for linear functional regression.
See e.g.~the upper bound on $\mathcal S_2(\lambda,z)$ in \citep{caponnetto2007optimal}, the second inequality in the proof of Theorem 5 in \citep{smale2007learning}, or Lemma 2 in \citep{cai2012minimax}.
Nevertheless, in our applications the choice of $\lambda_n$ is significantly smaller than conventional choices (for inference purposes), and when $\lambda_n = o(n^{-1})$ existing arguments cannot be directly applied,
without additional conditions such as Assumption (B3) for some $\kappa<1/2$. Therefore, we give the complete proof of Lemma \ref{lem:empirical-population-cov} below for our purposes.

\begin{proof}[Proof of Lemma \ref{lem:empirical-population-cov}.]
We use Hoeffding's inequality for bounded self-adjoint operators, as developed and stated in \citep{minsker2017some}.
For every $i\in[n]$, define
$$
W_i := (T+\lambda_n I)^{-1/2}Z_i = (T+\lambda_n I)^{-1/2}L_{K^{1/2}}Z_i.
$$
This gives rise to
\begin{align}
\frac{1}{n}\sum_{i=1}^n W_i\otimes W_i &= (T+\lambda_n I)^{-1/2}\hat T_n (T+\lambda_n I)^{-1/2};\label{eq:epcov-0}\\
\mathbb E[W_i\otimes W_i] &= (T+\lambda_n I)^{-1/2} T (T+\lambda_n I)^{-1/2}, \;\;\;\;\;\;\forall i\in[n]\label{eq:epcov-0half}.
\end{align}

By Assumption (B3), it holds almost surely that 
\begin{align}
\|W_i\otimes W_i\|_\op&=\left\|\sum_{j,k=1}^{\infty}\frac{\langle Z_i,\psi_j\rangle\langle Z_i,\psi_k\rangle}{\sqrt{s_j+\lambda_n}\sqrt{s_k+\lambda_n}}\psi_j\otimes\psi_k\right\|_\op\nonumber\\
&\leq \max_{j,k\geq 1}s_j^{-1/2}s_k^{-1/2}|\langle Z_i,\psi_j\rangle|\langle Z_i,\psi_k\rangle|\leq n^{2\kappa}.
\label{eq:proof-epcov-1}
\end{align}

Let 
$$
A_n := \sum_{i=1}^n\mathbb E[(W_i\otimes W_i)^2] =n\times\sum_{j,k=1}^{\infty}\frac{\mathbb E[\langle Z,Z\rangle\langle Z,\psi_j\rangle\langle Z,\psi_k\rangle]}{(s_j+\lambda_n)(s_k+\lambda_n)}\psi_j\otimes \psi_k.
$$
It is easy to verify that $\|A_n\|_\op \geq n\mathbb E[\langle Z,\psi_1\rangle^4]/(s_1+\lambda_n)=ns_1^2/(s_1+\lambda_n)\geq c_sn/2$ for sufficiently large $n$. 
On the other hand, Because $\langle Z,Z\rangle\leq 1$ almost surely thanks to Assumption (A3), it holds that
\begin{equation}
\tr(A_n) \leq n\times \sum_{j=1}^{\infty} \frac{s_j}{(s_j+\lambda_n)^2} \lesssim n^{(2r+1)/r},
\label{eq:proof-epcov-2}
\end{equation}
where the last inequality holds because $s_j\asymp j^{-2r}$ thanks to Assumption (A1), and $\lambda_n\gtrsim n^{-2}$.
Invoking \citep[Theorem 3.1]{minsker2017some}, it holds with probability $1-n^{-1}$ that
\begin{align*}
\left\|\frac{1}{n}\sum_{i=1}^n W_i\otimes W_i-\mathbb E[W_i\otimes W_i]\right\|_\op &\lesssim n^{-\frac{1-2\kappa}{2}}\log n = o(1). 
\end{align*}
Subsequently, incorporating Eqs.~(\ref{eq:epcov-0},\ref{eq:epcov-0half}), the above inequality implies
$$
\left\|(T+\lambda_n)^{-1/2}(T-\hat T_n)(T+\lambda_n^{-1/2})\right\|_\op = o(1).
$$
This proves Lemma \ref{lem:empirical-population-cov}, by selecting $n_0$ sufficiently large such that the right-hand side of the above inequality is upper bounded by $1/2$.
\end{proof}

\begin{lemma}
Assume Assumption (A1) holds and $\lambda_n$ are chosen such that Eq.~(\ref{eq:lambdan-1}) holds. Let $\varepsilon_i = Y_i - \langle X_i,\beta_0\rangle$. If \emph{either} one of the following conditions are satisfied:
\begin{enumerate}
\item $\{\varepsilon_i\}_{i=1}^n$ are Normally distributed, or
\item Assumption (B3) holds for some $\kappa<1/4$;
\end{enumerate}
then for every $x\in\supp(P_X)$ and $z=L_{K^{1/2}}x$, as $n\to\infty$, 
$$
\frac{1}{\tilde\sigma_n(z)\sqrt{n}}\sum_{i=1}^n \varepsilon_i g(z,Z_i)\;\;\overset{d}{\to}\;\; \mathcal N(0,1),
$$
where $Z_i=L_{K^{1/2}}X_i$, $g(z,Z)=\langle z,(\hat T_n+\lambda_n I)^{-1}Z\rangle$ and $\tilde\sigma_n(z)^2 = n^{-1}\sum_{i=1}^n \sigma_\varepsilon^2g(z,Z_i)^2$.
\label{lem:clt}
\end{lemma}

\begin{proof}[Proof of Lemma \ref{lem:clt}.]
It is easy to see that $\sum_{i=1}^n\varepsilon_i g(z,Z_i)$ is the sum of $n$ independent, centered random variables with $\varsigma_n^2 = \mathbb E[(\sum_{i=1}^n\varepsilon_ig(z,Z_i))^2] = n\tilde\sigma_n(z)^2$.
If $\{\varepsilon_i\}_{i=1}^n$ are independently distributed Normal random variables then the lemma obviously holds.

When $\{\varepsilon_i\}_{i=1}^n$ are not necessarily Normally distributed, 
we will verify Linderberg's condition for central limit theorem of independently but not identically distributed random variables.

If not explicitly mentioned, all probabilities and expectations are conditioned on $\{Z_i\}_{i=1}^n$ in the rest of this proof.
Fix arbitrary $\chi>0$. Because $\frac{1}{n}\sum_{i=1}^ng(z,Z_i)^2 = \tilde\sigma_n(z)^2/\sigma_\varepsilon$,
we know that for all but $m\leq 2/(\chi^2\sigma_\varepsilon^2)$ samples (denoted as $\mI\subseteq[n]$ such that $|\mI|=m$), 
$$
\big|g(z,Z_i)\big| \leq \frac{1}{2}\chi \sqrt{n}\tilde\sigma_n(z),\;\;\;\;\;\;\forall i\notin\mI.
$$
Because $z\neq 0$, Lemma \ref{lem:empirical-population-cov} implies that $\tilde\sigma_n(z)\overset{p}{\to} \langle z, T^{-1} z\rangle >0$ and therefore, applying Assumption (A1), it holds 
with probability $1-O_P(n^{-1})$ over the randomness of $\{Z_i\}_{i=1}^n$ that
\begin{equation}
\Pr\left[\big|\varepsilon_i g(z,Z_i)\big| \geq \chi \varsigma_n\right] \lesssim \exp\{-\Omega(n)\},\;\;\;\;\;\;\forall i\notin \mI
\label{eq:proof-clt-1}
\end{equation}
where in the $\Omega(\cdot)$ notation we hide dependency on $\chi$ and $\langle z,T^{-1}z\rangle$, both strictly positive constants independent of $n$ and $\{Z_i\}$.

{\color{red}
On the other hand, Assumptions (A3), (B3) and Lemma \ref{lem:empirical-population-cov} imply that, with probability $1-O_P(T^{-1})$ over the randomness of $\{Z_i\}_{i=1}^n$, 
\begin{equation}
g(z,Z_i)^2 = O_P(n^{4\kappa}),\;\;\;\;\;\;\forall i\in[n].
\label{eq:proof-clt-2}
\end{equation}
Combining Eqs.~(\ref{eq:proof-clt-1},\ref{eq:proof-clt-2}), we obtain with probability $1-O_P(n^{-1})$ over the randomness of $\{Z_i\}_{i=1}^n$ that
\begin{align*}
\frac{1}{\varsigma_n^2}\sum_{i=1}^n \mathbb E\left[(\varepsilon_ig(z,Z_i)^2\big| |\varepsilon_ig(z,Z_i)|>\chi\varsigma_n\right] &\leq \frac{1}{n\tilde\sigma_n(z)^2}\left[\frac{2}{\chi^2\sigma_\varepsilon^2}e^{-\Omega(n)} + O(n^{4\kappa}) \right]\to 0
\end{align*}
as $n\to\infty$. This verifies the Lindeberg's condition for central limit theorem.
}
\end{proof}


\subsection{Proof of Theorem \ref{thm:simple}}

We are now ready to prove Theorem \ref{thm:simple}.
With $Y_i=\langle X_i,\beta_0\rangle+\varepsilon_i = \langle Z_i,\theta_0\rangle+\varepsilon_i$, it holds that
$$
\hat\theta_n = (\hat T_n+\lambda_n I)^{-1}\hat T_n\theta_0 + \frac{1}{n}\sum_{i=1}^n\varepsilon_i (\hat T_n+\lambda_n I)^{-1}Z_i.
$$
Subsequently,
\begin{align}
\hat\theta_n-\theta_0 &= \underbrace{-\lambda_n (\hat T_n+\lambda_n I)^{-1}\theta_0}_{=: b_n} + \underbrace{\frac{1}{n}\sum_{i=1}^n\varepsilon_i(\hat T_n+\lambda_n I)^{-1}Z_i}_{=:v_n}.
\end{align}

For any test function $x\in\supp(P_X)$, let $z=L_{K^{1/2}}x=\sum_{j=1}^{n}\hat\nu_j\hat\psi_j$, where $\hat\nu_j=\langle z,\hat\psi_j\rangle$.
Let also $\hat\vartheta_j = \langle\theta_0,\hat\psi_j\rangle$.
Slightly abusing notations by denoting $b_n(z) := \langle b_n,z\rangle = \int_I b_n(t)z(t)\ud t$, we have that
\begin{align}
\big|b_n(z)\big| &= \lambda_n\left|\left\langle(\hat T_n+\lambda_n I)^{-1}\sum_{j=1}^n\hat\vartheta_j\hat\psi_j, \sum_{j=1}^n\hat\nu_j\hat\psi_j\right\rangle\right|
= \lambda_n \left|\sum_{j=1}^n \frac{\hat\nu_j\hat\vartheta_j}{\hat s_j+\lambda_n}\right|\label{eq:proof-simple-1}\\
&\leq \lambda_n \sqrt{\sum_{j=1}^n\hat\vartheta_j^2}\sqrt{\sum_{j=1}^n\frac{\hat\nu_j^2}{(\hat s_j+\lambda_n)^2}}\label{eq:proof-simple-2}\\
&\leq \sqrt{\lambda_n}\cdot \sqrt{\langle z,(\hat T_n+\lambda_n I)^{-1}z\rangle} \leq ??.\label{eq:proof-simple-3}
\end{align}
In the above derivations, Eq.~(\ref{eq:proof-simple-1}) holds because $\{\hat\psi_j\}_{j=1}^n$ is orthonormal with respect to $L_2(I)$ and $\hat T_n=\sum_{j=1}^n\hat s_j\hat\psi_j\otimes\hat\psi_j$;
Eq.~(\ref{eq:proof-simple-2}) holds by the Cauchy-Schwarz inequality;
Eq.~(\ref{eq:proof-simple-3}) holds by noting that $\lambda_n/(\sqrt{\hat s_j+\lambda_n})\leq\sqrt{\lambda_n}$ for all $j$, and furthermore
$$
\langle z, (\hat T_n+\lambda_n I)^{-1}z\rangle = \left\langle\sum_{j=1}^n\hat\nu_j\hat\psi_j, \sum_{j=1}^n\frac{\hat\nu_j}{\hat s_j+\lambda_n}\hat\psi_j\right\rangle = \sum_{j=1}^n\frac{\hat\nu_j^2}{\hat s_j+\lambda_n}.
$$

We next analyze the term involving $v_n$ on an arbitrary test function $x\in\supp(P_X)$ with $z=L_{K^{1/2}}x=\sum_{j=1}^n\hat\nu_j\hat\psi_j$, where $\hat\nu_j=\langle z,\hat\psi_j\rangle$.
Let $v_n(z)=\langle v_n,z\rangle = \int_I v_n(t)z(t)\ud t$. We have that
$$
v_n(z) = \frac{1}{n}\sum_{i=1}^n \varepsilon_i \langle z, (\hat T_n+\lambda_n I)^{-1}Z_i\rangle.
$$
Note that $v_n(z)$ is \emph{pivotal}, meaning that it does not depend on the unknown $\beta_0\in\Theta_2(\mH)$ or $\theta_0=L_{K^{1/2}}\beta_0$.
It should also be noted that $\langle z,(\hat T_n+\lambda_n I)^{-1}Z_i\rangle$ are independent of $\{\varepsilon_i\}_{i=1}^n$.
{\color{red}Using central limit theorem with Lyapnov conditions}, it holds that 
\begin{equation}
\sqrt{n}\hat\sigma_n(z)^{-1}v_n(z) \;\;\overset{d}{\to}\;\; \mathcal N(0,1)
\label{eq:proof-simple-4}
\end{equation}
where
$$
\hat\sigma_n(z)^2 = \frac{1}{n}\sum_{i=1}^n \langle z,(\hat T_n+\lambda_n I)^{-1}Z_i\rangle^2 = \langle z, (\hat T_n+\lambda_n I)^{-1}\hat T_n(\hat T_n+\lambda_n I)^{-1} z\rangle.
$$
Using Slutsky's theorem and the condition that $\lambda_n n\to 0$, Eq.~(\ref{eq:proof-simple-4}) implies
$$
\sqrt{n}\hat\sigma_n(z)^{-1}\left(b_n(z)+v_n(z)\right)\;\;\overset{d}{\to}\;\;\mathcal N(0,1).
$$
Because the above convergence in distribution holds for any $z$ and $\beta_0\in\Theta_2(\mH)$, we have proved the honesty of our constructed confidence intervals as stipulated in Eq.~(\ref{eq:simple-validity}) of Theorem \ref{thm:simple}.

\bibliographystyle{apalike}
\bibliography{refs}

\end{document}